\paragraph{Limitations.} Ideally, steering should be targeted; shifting from a “yes” to a “no” should affect only questions initially answered with “yes”. While this is often achieved, it is not perfect. We also notice a tradeoff between the steering strength and the quality of generated responses. 
This approach focuses on changing the model answers or style, however, it can be extended to other applications, such as addressing biases or mitigating safety concerns. It is also interesting to extend our study to larger and more recent MLLMs, eventually with different architectures.
\paragraph{Conclusion.} Gaining insights into the mechanisms of recent foundation models is crucial for advancing AI research. In this study, we delve into how MLLMs’ internal representations evolve during fine-tuning, demonstrating that concepts learned post-fine-tuning can often be recovered from the original model by navigating the feature space. This approach inspired our final investigation, where we steer model behavior by modifying features directly, without additional training. We show that this could be used to change model answers, or have richer captions focusing on different aspects of the image. We hope this work motivates further research into understanding these models and developing more efficient methods for adapting or refining them, moving beyond traditional training practices.





